\Needspace{11\baselineskip}
\begin{center}
\textit{At least half of all assignable characters for a given determinant cannot correspond to properly primitive genera (positive for a negative determinant).}
\end{center}

\Needspace{8\baselineskip}
\Article{261}
\textsc{Theorem.} \textit{For a positive nonsquare determinant, at least half of all assignable characters cannot correspond to any properly primitive genera; for a negative determinant, to any positive properly primitive genera.}

\textit{Proof.} Let $m$ be the number of all properly primitive (positive) genera of determinant $D$, and $k$ the number of classes in each genus, so that $km$ is the number of all properly primitive (positive) classes; let $n$ be the number of all distinct assignable characters for this determinant. Then
\SourceLeaf{297}
by Article 258 the number of all properly primitive (positive) ambiguous classes is $\frac12 n$. Hence by the preceding article the number of all properly primitive classes obtainable by doubling such a class is $\frac{2km}{n}$. But by Article 247 all these classes belong to the principal genus, which contains $k$ classes. Thus, if all classes of the principal genus can be obtained by doubling some class (as will be proved below to hold in every case), we have $\frac{2km}{n}=k$, or $m=\frac12 n$; but certainly we cannot have $\frac{2km}{n}>k$, and therefore not $m>\frac12 n$. Since, then, the number of all properly primitive (positive) genera is certainly no greater than half the number of all assignable characters, at least half of these characters cannot correspond to such genera. \textit{This proves the theorem.} --- It must be kept in mind, however, that this does not yet show that half of all assignable characters actually correspond to properly primitive (positive) genera. The truth of this very important proposition can only be brought out below from the deepest mysteries of numbers.

Since for a negative determinant there are always as many negative as positive genera, clearly no more than half of all assignable characters can belong to negative properly primitive genera. We will discuss this below, as well as improperly primitive genera. Finally, the theorem does not extend to positive square determinants: in this case it is easy to see that genera actually correspond to every assignable character.

\Needspace{11\baselineskip}
\begin{center}
\textit{A second proof of the fundamental theorem and the other theorems concerning the residues $-1$, $+2$, $-2$.}
\end{center}

\Needspace{8\baselineskip}
\Article{262}
Thus, when only two distinct characters can be assigned for a given nonsquare determinant $D$, only one corresponds to a properly primitive (positive) genus (which can only be the principal genus); the other belongs to no properly primitive (positive) form of that determinant. This happens for the determinants $-1$, $2$, $-2$, $-4$. prime numbers of the form $4n+1$ taken positively and those of the form $4n+3$ taken negatively; finally, for all odd powers of primes of the form $4n+1$ taken positively, and for powers of primes of the form $4n+3$ taken positively or negatively according as their exponents are even or odd. From this principle we can draw a new method of proving not only the fundamental theorem but also the other theorems of the preceding section concerning the residues $-1$, $+2$, $-2$. It is entirely different from the methods used in the
\SourceLeaf{298}
preceding section and seems no less elegant. We will pass over the determinant $-4$ and those that are powers of primes, since they teach us nothing new.

Thus for the determinant $-1$ there is no positive form with character $3,4$; for the determinant $+2$ there is no form at all with character $3\ \mathit{and}\ 5,8$; for the determinant $-2$ no positive form has character $5\ \mathit{and}\ 7,8$; for the determinant $+p$, if $p$ is a prime of the form $4n+1$, or for the determinant $-p$, if $p$ is a prime of the form $4n+3$, no properly primitive form (positive in the latter case) has character $Np$. We therefore prove the theorems of the preceding section as follows.

I. $-1$ is a nonresidue of every (positive) number of the form $4n+3$. For if $-1$ were a residue of such a number $A$, setting $-1=BB-AC$ would give $(A,B,C)$ as a positive form of determinant $-1$ with character $3,4$.

II. $-1$ is a residue of every prime $p$ of the form $4n+1$. For the character of the form $(-1,0,p)$, like that of every properly primitive form of determinant $p$. is $Rp$, and therefore $-1Rp$.

III. Both $+2$ and $-2$ are residues of every prime $p$ of the form $8n+1$. For either the forms $\left(8,\allowbreak 1,\allowbreak \frac{1-p}{8}\right)$, $\left(-8,\allowbreak 1,\allowbreak \frac{p-1}{8}\right)$ or the forms $\left(8,\allowbreak 3,\allowbreak \frac{9-p}{8}\right)$, $\left(-8,\allowbreak 3,\allowbreak \frac{p-9}{8}\right)$ are properly primitive (according as $n$ is odd or even), so their character is $Rp$. Thus $+8Rp$ and $-8Rp$, which also gives $2Rp$, $-2Rp$.

IV. $+2$ is a nonresidue of every number of the form $8n+3$ or $8n+5$. For if it were a residue of such a number $A$, there would be a form $(A,B,C)$ of determinant $+2$ with character $3\ \mathit{and}\ 5,8$.

V. Likewise $-2$ is a nonresidue of every number of the form $8n+5$ or $8n+7$; otherwise there would be a form $(A,B,C)$ of determinant $-2$ with character $5\ \mathit{and}\ 7,8$.

VI. $-2$ is a residue of every prime $p$ of the form $8n+3$. This proposition can be proved in two ways. \textit{First}, since by IV we have $+2Np$, and by I $-1Np$, we must have $-2Rp$. The \textit{second} proof comes from considering the determinant $+2p$, for which four characters can be assigned, namely $Rp$, $1\ \mathit{and}\ 3,8$; $Rp$, $5\ \mathit{and}\ 7,8$; $Np$, $1\ \mathit{and}\ 3,8$; $Np$, $5\ \mathit{and}\ 7,8$. Thus at least two cannot correspond to any genera. Now the form
\SourceLeaf{299}
$(1,0,-2p)$ has the first character, and the form $(-1,0,2p)$ the fourth, so the second and third must be rejected. Since the character of the form $(p,0,-2)$ relative to the number $8$ is $1\ \mathit{and}\ 3,8$, its character relative to $p$ can only be $Rp$, giving $-2Rp$.

VII. $+2$ is a residue of every prime $p$ of the form $8n+7$, which can be proved in two ways. \textit{First}, since I and V give $-1Np$, $-2Np$, we have $+2Rp$. \textit{Second}, since either $\left(8,\allowbreak 1,\allowbreak \frac{1+p}{8}\right)$ or $\left(8,\allowbreak 3,\allowbreak \frac{9+p}{8}\right)$ is a properly primitive form of determinant $-p$ (according as $n$ is even or odd), its character is $Rp$, and therefore $8Rp$ and $2Rp$.

VIII. Every prime $p$ of the form $4n+1$ is a nonresidue of every odd number $q$ that is a nonresidue of $p$. For if $p$ were a residue of $q$, clearly there would be a properly primitive form of determinant $p$ with character $Np$.

IX. Likewise, if any odd number $q$ is a nonresidue of a prime $p$ of the form $4n+3$, $-p$ is a nonresidue of $q$; otherwise there would be a positive properly primitive form of determinant $-p$ with character $Np$.

X. Every prime $p$ of the form $4n+1$ is a residue of every other prime $q$ that is a residue of $p$. If $q$ too is of the form $4n+1$, this follows at once from VIII; but if $q$ is of the form $4n+3$, $-q$ too is a residue of $p$ (by II), and therefore $pRq$ (by IX).

XI. If any prime $q$ is a residue of another prime $p$ of the form $4n+3$, $-p$ is a residue of $q$. For if $q$ is of the form $4n+1$; VIII gives $pRq$, and therefore (by II) $-pRq$. But the case where $q$ too is of the form $4n+3$ escapes this method, though it is easily settled by considering the determinant $+pq$. Of the four assignable characters for this determinant, $Rp$. $Rq$; $Rp$, $Nq$; $Np$, $Rq$; $Np$, $Nq$, two cannot correspond to any genera. Since the forms $(1,0,-pq)$, $(-1,0,pq)$ have respectively the first and fourth characters, the second and third cannot belong to any properly primitive form of determinant $pq$. Thus, since by hypothesis the character of the form $(q,0,-p)$ relative to the number $p$ is $Rp$, its character relative to $q$ must be $Rq$, and therefore $-pRq$. \textit{This proves the theorem.}

If in propositions VIII and IX $q$ is taken to denote a prime, those propositions together with X and XI give the fundamental theorem of the preceding section.
\Needspace{11\baselineskip}
\SourceLeaf{300}
\begin{center}
\textit{A closer determination of the half of the characters to which no genera can correspond.}
\end{center}

\Needspace{8\baselineskip}
\Article{263}
Having established the fundamental theorem by a new proof, we will show how to distinguish, for any given nonsquare determinant, the half of the characters that cannot correspond to any properly primitive (positive) forms. We can do this more briefly since its basis is already contained in the investigation of Articles 147--150. Let $ee$ be the greatest square dividing the proposed determinant $D$, and $D=D'ee$, so that $D'$ contains no square factor. Further, let $a$, $b$, $c$ etc. be all the odd prime divisors of $D'$, so that $D'$, apart from its sign, is either their product or twice that product. Let $\Omega$ denote the collection of individual characters $Na$, $Nb$, $Nc$ etc. alone when $D'\equiv1\ (\mathrm{mod.}\ 4)$; with the character $3,4$ adjoined when $D'\equiv3$ and $e$ is odd or twice an odd number; with $3,8$ and $7,8$ adjoined when $D'\equiv3$ and $e$ is divisible by four; with either the character $3\ \mathit{and}\ 5,8$ or the two $3,8$ and $5,8$ adjoined when $D'\equiv2\ (\mathrm{mod.}\ 8)$, according as $e$ is odd or even; finally, with either the character $5\ \mathit{and}\ 7,8$ or the two $5,8$ and $7,8$ adjoined when $D'\equiv6\ (\mathrm{mod.}\ 8)$, according as $e$ is odd or even. Then no complete character containing an odd number of the individual characters in $\Omega$ can correspond to properly primitive (positive) genera of determinant $D$. In all cases, individual characters expressing a relation to prime divisors of $D$ that do not divide $D'$ contribute nothing to the possibility or impossibility of genera. --- From the theory of combinations it is very easy to see that this indeed excludes half of all assignable complete characters.

These rules are proved as follows. From the principles of the preceding section, or the theorems proved again in the preceding article, it follows without difficulty that if $p$ is a prime (odd and positive) not dividing $D$ and having one of the rejected characters, $D'$ contains an odd number of factors that are nonresidues of $p$, so $D'$, and therefore also $D$, is a nonresidue of $p$. It is also easy to see that a product of any number of odd numbers relatively prime to $D$, none of which has any rejected character, cannot have such a character either. Conversely, any positive odd number relatively prime to $D$ that has one of the \SourceLeaf{301}rejected characters must therefore contain a prime factor with the same property, so $D$ is a nonresidue of that number. Thus, if there were a properly primitive (positive) form of determinant $D$ corresponding to a rejected character, $D$ would be a nonresidue of every positive odd number relatively prime to it represented by such a form, which clearly cannot be reconciled with the theorem of Article 154.

For examples, compare the classifications in Articles 231, 232; anyone can increase their number at will.
